\documentclass[12pt,final]{report}

% PAQUETES 

\usepackage{graphicx}
\setkeys{Gin}{draft=false}
\graphicspath{{./}}


\usepackage{amsmath}
\usepackage{amssymb}
\usepackage[spanish]{babel}
\usepackage[utf8]{inputenc}
\usepackage{lmodern}
\usepackage[T1]{fontenc}
\usepackage[margin=2.5cm]{geometry}
\usepackage{siunitx}
\usepackage{setspace}
\usepackage{microtype}
\usepackage{amsthm}
\usepackage{hyperref}
\usepackage{fancyhdr}
\usepackage{tocloft}
\usepackage{enumitem}


\renewcommand{\cftsecleader}{\cftdotfill{\cftdotsep}}
\setlength{\cftbeforesecskip}{6pt}
\renewcommand{\contentsname}{\Huge \textbf{Índice}}
\renewcommand{\cftchapfont}{\bfseries}
\renewcommand{\cftsecfont}{}
\renewcommand{\cftchappagefont}{\bfseries}

\setlength{\cftbeforechapskip}{10pt}
\setlength{\cftbeforesecskip}{4pt}

\setlength{\parindent}{0pt}
\setlength{\parskip}{6pt}



% ENTORNOS MATEMATICOS

\theoremstyle{plain}
\newtheorem{theorem}{Teorema}[chapter]
\newtheorem{proposition}[theorem]{Proposición}
\newtheorem{lemma}[theorem]{Lema}

\theoremstyle{definition}
\newtheorem{definition}[theorem]{Definición}

% DATOS 

\title{Demostraciones Matemáticas}
\author{Álvaro Yamil Manzo Ayora}
\date{Marzo 2026}

\begin{document}
\hypersetup{pageanchor=false}

% PORTADA 


\begin{titlepage}
\begin{center}

% LOGOS

\includegraphics[width=3.5cm]{logo_unam_negro (1).png}
\hspace{1.5cm}
\includegraphics[width=3.5cm]{logo_enp8.jpeg}

\vspace{1cm}

{\Large \textbf{Universidad Nacional Autónoma de México}}

\vspace{0.3cm}

{\large Escuela Nacional Preparatoria No. 8}

\vspace{0.3cm}

{\large ``Miguel E. Schulz''}

\vspace{2cm}

{\Huge \textbf{Demostraciones Matemáticas}}

\vspace{0.5cm}
\rule{\textwidth}{0.6pt}
\vspace{0.5cm}

\vspace{2cm}

\large
\textbf{Autor:} Álvaro Yamil Manzo Ayora

\vfill

{\large Ciudad de México, México}

\vspace{0.3cm}

{\large Marzo 2026}

\end{center}
\end{titlepage}

\pagenumbering{gobble}

\begin{center}
\setcounter{tocdepth}{1}
\tableofcontents
\vspace{1cm}
\end{center}

\clearpage
\pagenumbering{arabic}
\hypersetup{pageanchor=true}


% PROLOGO 

\chapter*{Prólogo}
\addcontentsline{toc}{chapter}{Prólogo}

En este documento (codificado en \LaTeX), creado el 27 de marzo de 2026 por Álvaro Yamil Manzo Ayora, estudiante de bachillerato, evidenciaré mis avances en la lógica matemática.

Aunque las demostraciones presentadas hasta el momento sean bastante básicas y algunas puedan parecer evidentes, considero importante mostrarlas desde el inicio para que, en su debido momento, pueda mirar hacia atrás y recordar lo que fui y lo que soy mediante el progreso reflejado en este archivo.

Igualmente, espero que las personas que no estén familiarizadas con estos temas puedan avanzar en las demostraciones con mayor facilidad y comenzar a desarrollar el pensamiento abstracto.

Actualmente este documento se encuentra en código abierto en mi repositorio de GitHub:

\begin{center}

\url{https://github.com/Alvaro-Manzo/DEMOSTRACIONES}

\end{center}

Donde se encuentra el código de este archivo también existen demás códigos en LaTeX que podrían ser de interés para cualquier persona curiosa; le recomiendo estar al tanto del repositorio.

Espero continuar contribuyendo a este proyecto compartiendo mis pequeños conocimientos y, en algún momento, poder compartirlo con todos, ayudando a las demás generaciones y tratando de hacer la matemática más accesible para quienes la consideran complicada o aburrida.

No te desesperes si algo no se entiende al principio. Todos, incluso los grandes matemáticos, alguna vez no entendimos nada. Con paciencia y práctica, el pensamiento matemático se vuelve cada vez más claro.

\vspace{1.5em}

\begin{center}

\textbf{Le recomiendo mantener presente la siguiente frase mientras lee este documento:}

\vspace{1em}

\textit{``La belleza de las matemáticas solo se muestra a sus seguidores más pacientes.''}

\textsc{Maryam Mirzakhani}

\end{center}

\newpage

% CAPITULO 1

\chapter{Demostraciones aritméticas}

\section{Propiedades Básicas}

\subsection{Demostración de $1 = 0.999\ldots$}

\begin{proposition}
El número $1$ es igual al número decimal periódico $0.999\ldots$.
\end{proposition}

\begin{proof}

Sabemos que

\[
\frac{1}{3} = 0.333\ldots
\]

Multiplicando ambos lados por $3$:

\[
3 \cdot \frac{1}{3} = 3 \cdot 0.333\ldots
\]

\[
\frac{3}{3} = 0.999\ldots
\]

Simplificando la fracción:

\[
1 = 0.999\ldots
\]

\end{proof}

\newpage

\subsection{Demostración de $1>0$}

\begin{proposition}
En los números reales, se cumple que $1>0$.
\end{proposition}

\begin{proof}

Observemos que

\[
1 = 1^2
\]

En los números reales, el cuadrado de todo número distinto de cero es positivo. Como $1 \neq 0$, se sigue que

\[
1^2 > 0
\]

Por lo tanto,

\[
1 > 0
\]

\end{proof}

\newpage

\subsection{Demostración de $0\cdot a = 0$}

\begin{proposition}
Para cualquier número $a$ se cumple que
\[
0\cdot a = 0.
\]
\end{proposition}

\begin{proof}

Sabemos que

\[
0 = 0+0
\]

Multiplicamos ambos lados por $a$:

\[
a\cdot0 = a(0+0)
\]

Aplicamos la propiedad distributiva:

\[
a(0+0) = a\cdot0 + a\cdot0
\]

Entonces

\[
a\cdot0 = a\cdot0 + a\cdot0
\]

Restamos $a\cdot0$ en ambos lados:

\[
a\cdot0 - a\cdot0 = a\cdot0 + a\cdot0 - a\cdot0
\]

Simplificando:

\[
0 = a\cdot0
\]

\end{proof}

\newpage

\subsection{Demostración de que $\sqrt{2}$ es irracional}

\begin{proposition}
El número $\sqrt{2}$ es irracional.
\end{proposition}

\begin{proof}

Procedemos por contradicción.

Supongamos que $\sqrt{2}$ es racional.

Entonces puede escribirse como

\[
\sqrt{2} = \frac{a}{b}
\]

donde $a$ y $b$ son enteros sin factores comunes.

Elevando al cuadrado:

\[
2 = \frac{a^2}{b^2}
\]

Multiplicando por $b^2$:

\[
2b^2 = a^2
\]

Esto implica que $a^2$ es par, por lo tanto $a$ es par.

Entonces existe un entero $k$ tal que

\[
a = 2k
\]

Sustituyendo:

\[
2b^2 = 4k^2
\]

\[
b^2 = 2k^2
\]

Esto implica que $b$ también es par.

Pero entonces $a$ y $b$ son ambos pares, lo cual contradice que la fracción fuera irreducible.

Por lo tanto

\[
\sqrt{2} \text{ es irracional}
\]

\end{proof}

\newpage
\subsection{Demostración de que la suma de dos números pares es par}

\begin{proposition}
La suma de dos números pares es par.
\end{proposition}

\begin{proof}

Sea $a$ y $b$ dos números pares. Por definición, existen enteros $k$ y $m$ tales que:

\[
a = 2k, \quad b = 2m
\]

Entonces:

\[
a + b = 2k + 2m
\]

Agrupamos:

\[
a + b = 2(k + m)
\]

Sabemos que la suma de dos enteros es un entero, es decir, $k + m \in \mathbb{Z}$.

Por lo tanto, $a+b$ es múltiplo de 2.

Esto implica que:

\[
a + b \text{ es par}
\]

\end{proof}

\newpage

\subsection{Demostración de que la suma de dos números impares es par}

\begin{proposition}
La suma de dos números impares es par.
\end{proposition}

\begin{proof}

Sean $a$ y $b$ números impares. Entonces:

\[
a = 2k + 1, \quad b = 2m + 1
\]

Sumamos:

\[
a + b = (2k + 1) + (2m + 1)
\]

\[
a + b = 2k + 2m + 2
\]

Factorizamos:

\[
a + b = 2(k + m + 1)
\]

Como $k+m+1$ es un entero, el resultado es múltiplo de 2.

Por lo tanto, la suma es par.

\end{proof}

\newpage

\subsection{Demostración de que el cuadrado de un número impar es impar}

\begin{proposition}
El cuadrado de un número impar es impar.
\end{proposition}

\begin{proof}

Sea $n$ impar. Entonces:

\[
n = 2k + 1
\]

Elevamos al cuadrado:

\[
n^2 = (2k + 1)^2
\]

Desarrollamos:

\[
n^2 = 4k^2 + 4k + 1
\]

Factorizamos:

\[
n^2 = 2(2k^2 + 2k) + 1
\]

Observamos que $2k^2 + 2k$ es entero.

Por lo tanto:

\[
n^2 = 2m + 1
\]

Esto prueba que $n^2$ es impar.

\end{proof}

\newpage

\subsection{Demostración de que si $ab=0$, entonces $a=0$ o $b=0$}

\begin{proposition}
Si $a$ y $b$ son números reales y $ab=0$, entonces $a=0$ o $b=0$.
\end{proposition}

\begin{proof}

Sea $ab=0$.

Si $a=0$, entonces ya se cumple la conclusión.

Supongamos ahora que $a \neq 0$. Como estamos trabajando en los números reales, podemos dividir entre $a$. Entonces:

\[
\frac{ab}{a} = \frac{0}{a}
\]

De aquí obtenemos

\[
b = 0
\]

Por lo tanto, si $ab=0$, necesariamente se cumple que

\[
a=0 \quad \text{o} \quad b=0
\]

\end{proof}

\newpage

\subsection{Demostración de la suma de los primeros $n$ números naturales}

\begin{proposition}
\[
1+2+3+\cdots+n = \frac{n(n+1)}{2}
\]
\end{proposition}

\begin{proof}

Usaremos inducción matemática.

\textbf{Caso base:}

Para $n=1$:

\[
1 = \frac{1(1+1)}{2} = 1
\]

Se cumple.

\textbf{Paso inductivo:}

Supongamos que:

\[
1+2+\cdots+n = \frac{n(n+1)}{2}
\]

Entonces:

\[
1+2+\cdots+n+(n+1)
\]

Sustituimos:

\[
= \frac{n(n+1)}{2} + (n+1)
\]

Factorizamos:

\[
= \frac{(n+1)(n+2)}{2}
\]

Por lo tanto, la fórmula se cumple para $n+1$.

\end{proof}

\newpage

\subsection{Demostración de que no existe el mayor número natural}

\begin{proposition}
No existe un número natural máximo.
\end{proposition}

\begin{proof}

Supongamos que existe un número natural máximo $N$.

Entonces ningún número natural puede ser mayor que $N$.

Pero consideremos el número:

\[
N+1
\]

Claramente:

\[
N+1 > N
\]

Esto contradice la hipótesis de que $N$ es el mayor número natural.

Por lo tanto, no existe tal número.

\end{proof}

\newpage


\subsection{Demostración de que $\sqrt{3}$ es irracional}

\begin{proposition}
El número $\sqrt{3}$ es irracional.
\end{proposition}

\begin{proof}

Procedemos por contradicción.

Supongamos que $\sqrt{3}$ es racional. Entonces existen enteros $a$ y $b$ sin factores comunes tales que:

\[
\sqrt{3} = \frac{a}{b}
\]

Elevamos al cuadrado ambos lados:

\[
3 = \frac{a^2}{b^2}
\]

Multiplicamos por $b^2$:

\[
3b^2 = a^2
\]

Esto implica que $a^2$ es múltiplo de $3$. Por lo tanto, $a$ es múltiplo de $3$ (si un número no es múltiplo de $3$, su cuadrado deja residuo $1$ al dividirse entre $3$). Entonces existe $k \in \mathbb{Z}$ tal que:

\[
a = 3k
\]

Sustituimos en la ecuación $3b^2 = a^2$:

\[
3b^2 = (3k)^2 = 9k^2
\]

Dividimos entre $3$:

\[
b^2 = 3k^2
\]

Esto implica que $b^2$ es múltiplo de $3$, y por lo tanto $b$ también es múltiplo de $3$.

Hemos llegado a una contradicción, pues $a$ y $b$ son ambos múltiplos de $3$, lo que contradice la hipótesis de que la fracción $\frac{a}{b}$ es irreducible.

Por lo tanto, $\sqrt{3}$ es irracional.

\end{proof}

\newpage

\subsection{Demostración de que el producto de dos números impares es impar}

\begin{proposition}
El producto de dos números impares es impar.
\end{proposition}

\begin{proof}

Sean $a$ y $b$ números impares. Por definición, existen enteros $k$ y $m$ tales que:

\[
a = 2k + 1, \qquad b = 2m + 1
\]

Calculamos su producto:

\[
a \cdot b = (2k + 1)(2m + 1)
\]

Desarrollamos usando la propiedad distributiva:

\[
a \cdot b = 4km + 2k + 2m + 1
\]

Factorizamos $2$ de los primeros tres términos:

\[
a \cdot b = 2(2km + k + m) + 1
\]

Sea $p = 2km + k + m$. Como $k$ y $m$ son enteros, $p$ es un entero. Entonces:

\[
a \cdot b = 2p + 1
\]

Esto es precisamente la forma de un número impar. Por lo tanto, el producto de dos números impares es impar.

\end{proof}

\newpage

\subsection{Demostración de que $(-a)(-b) = ab$}

\begin{proposition}
Para cualquier par de números reales $a$ y $b$, se cumple que:
\[
(-a)(-b) = ab
\]
\end{proposition}

\begin{proof}

Comenzamos con la expresión $(-a)(-b)$.

Sumamos y restamos $ab$ dentro de la expresión de manera inteligente. Primero, usamos la propiedad distributiva y el hecho de que $(-a) = (-1) \cdot a$ y $(-b) = (-1) \cdot b$:

\[
(-a)(-b) = [(-1) \cdot a] \cdot [(-1) \cdot b]
\]

Por la propiedad asociativa y conmutativa del producto:

\[
(-a)(-b) = (-1) \cdot (-1) \cdot a \cdot b
\]

Ahora demostramos que $(-1)(-1) = 1$. Sabemos que:

\[
0 = 0 \cdot (-1) = (1 + (-1)) \cdot (-1) = 1 \cdot (-1) + (-1) \cdot (-1) = -1 + (-1)(-1)
\]

Sumamos $1$ a ambos lados de la igualdad $-1 + (-1)(-1) = 0$:

\[
(-1)(-1) = 1
\]

Por lo tanto:

\[
(-a)(-b) = 1 \cdot a \cdot b = ab
\]

\end{proof}

\newpage

\subsection{Demostración de la desigualdad de las medias (AM-GM) para dos números}

\begin{proposition}
Para cualesquiera números reales no negativos $x$ e $y$, se cumple que:
\[
\frac{x+y}{2} \geq \sqrt{xy}
\]
con igualdad si y solo si $x = y$.
\end{proposition}

\begin{proof}

Observamos que para cualquier número real, su cuadrado es no negativo. En particular:

\[
(\sqrt{x} - \sqrt{y})^2 \geq 0
\]

Desarrollamos el cuadrado:

\[
(\sqrt{x})^2 - 2\sqrt{x}\sqrt{y} + (\sqrt{y})^2 \geq 0
\]

\[
x - 2\sqrt{xy} + y \geq 0
\]

Sumamos $2\sqrt{xy}$ a ambos lados:

\[
x + y \geq 2\sqrt{xy}
\]

Dividimos entre $2$:

\[
\frac{x+y}{2} \geq \sqrt{xy}
\]

La igualdad se da cuando $(\sqrt{x} - \sqrt{y})^2 = 0$, lo que ocurre si y solo si $\sqrt{x} = \sqrt{y}$, es decir, $x = y$.

\end{proof}

\newpage


\subsection{Demostración de que $\sqrt{ab} \leq \frac{a+b}{2}$ (equivalente a AM-GM)}

\begin{proposition}
Para $a, b \geq 0$:
\[
\sqrt{ab} \leq \frac{a+b}{2}
\]
\end{proposition}

\begin{proof}

Esta proposición es equivalente a la desigualdad AM-GM ya demostrada, pero presentamos una demostración alternativa directa.

Observamos que:

\[
\frac{a+b}{2} - \sqrt{ab} = \frac{a+b - 2\sqrt{ab}}{2} = \frac{(\sqrt{a})^2 - 2\sqrt{a}\sqrt{b} + (\sqrt{b})^2}{2}
\]

Reconocemos el numerador como un binomio al cuadrado:

\[
\frac{a+b}{2} - \sqrt{ab} = \frac{(\sqrt{a} - \sqrt{b})^2}{2}
\]

Como el cuadrado de cualquier número real es mayor o igual que cero:

\[
(\sqrt{a} - \sqrt{b})^2 \geq 0
\]

Por lo tanto:

\[
\frac{a+b}{2} - \sqrt{ab} \geq 0
\]

Sumamos $\sqrt{ab}$ a ambos lados:

\[
\frac{a+b}{2} \geq \sqrt{ab}
\]

La igualdad ocurre cuando $\sqrt{a} - \sqrt{b} = 0$, es decir, $a = b$.

\end{proof}

\newpage

\subsection{Demostración de que $n^2 \geq n$ para todo entero $n \geq 1$}

\begin{proposition}
Para todo número entero $n \geq 1$, se cumple que:
\[
n^2 \geq n
\]
\end{proposition}

\begin{proof}

Demostramos por inducción sobre $n$.

\textbf{Caso base:} Para $n = 1$:
\[
1^2 = 1 \geq 1
\]
Se cumple la igualdad.

\textbf{Paso inductivo:} Supongamos que para algún $k \geq 1$ se cumple $k^2 \geq k$. Demostremos que $(k+1)^2 \geq k+1$.

Desarrollamos $(k+1)^2$:

\[
(k+1)^2 = k^2 + 2k + 1
\]

Por hipótesis de inducción, $k^2 \geq k$, entonces:

\[
k^2 + 2k + 1 \geq k + 2k + 1 = 3k + 1
\]

Como $k \geq 1$, tenemos $3k + 1 \geq k + 1$ (ya que $3k + 1 - (k+1) = 2k \geq 0$). Por lo tanto:

\[
(k+1)^2 \geq k + 1
\]

Esto completa el paso inductivo.

Por el principio de inducción matemática, la proposición es válida para todo $n \geq 1$.

\end{proof}

\newpage

\subsection{Demostración de que $2^n > n$ para todo $n \in \mathbb{N}$}

\begin{proposition}
Para todo número natural $n$, se cumple que:
\[
2^n > n
\]
\end{proposition}

\begin{proof}

Procedemos por inducción matemática.

\textbf{Caso base:} Para $n = 1$:
\[
2^1 = 2 > 1
\]
Se cumple la desigualdad.

\textbf{Paso inductivo:} Supongamos que para algún $k \geq 1$ se cumple $2^k > k$. Demostremos que $2^{k+1} > k+1$.

Observamos que:

\[
2^{k+1} = 2 \cdot 2^k
\]

Por hipótesis de inducción, $2^k > k$, entonces:

\[
2 \cdot 2^k > 2k
\]

Ahora necesitamos demostrar que $2k \geq k+1$ para $k \geq 1$:
\[
2k - (k+1) = k - 1 \geq 0 \quad \text{para todo } k \geq 1
\]
Por lo tanto, $2k \geq k+1$.

Entonces:

\[
2^{k+1} > 2k \geq k+1
\]

Por lo tanto, $2^{k+1} > k+1$, lo que completa el paso inductivo.

Así, la desigualdad se cumple para todo $n \in \mathbb{N}$.

\end{proof}

\newpage

\subsection{Demostración de que existen infinitos números primos}

\begin{proposition}
Existen infinitos números primos.
\end{proposition}

\begin{proof}

Procedemos por contradicción.

Supongamos que solo existen finitos números primos:
\[
p_1, p_2, p_3, \dots, p_n
\]

Construimos el siguiente número:
\[
N = p_1 p_2 p_3 \cdots p_n + 1
\]

Este número $N$ es mayor que todos los primos de la lista.

Ahora analizamos si $N$ es primo o compuesto.

\textbf{Caso 1:} Si $N$ es primo, entonces no está en la lista, lo cual contradice que ya teníamos todos los primos.

\textbf{Caso 2:} Si $N$ es compuesto, entonces debe tener un divisor primo.

Pero ningún primo $p_i$ divide a $N$, porque:
\[
N \equiv 1 \pmod{p_i}
\]

Esto significa que al dividir $N$ entre cualquier $p_i$, el residuo es $1$.

Por lo tanto, $N$ tiene un divisor primo que no está en la lista.

En ambos casos obtenemos una contradicción.

\[
\therefore \text{existen infinitos números primos}
\]

\end{proof}

\newpage


\subsection{Demostración de que el producto de un número par por cualquier entero es par}

\begin{proposition}
Si $a$ es un número par y $b$ es un número entero, entonces $ab$ es par.
\end{proposition}

\begin{proof}

Como $a$ es par, por definición existe un entero $k$ tal que:

\[
a = 2k
\]

Multiplicamos ambos lados por $b$:

\[
ab = (2k)b
\]

Usando la asociatividad y conmutatividad del producto, obtenemos:

\[
ab = 2(kb)
\]

Como $k$ y $b$ son enteros, su producto $kb$ también es un entero. Por lo tanto, $ab$ tiene la forma $2m$, donde $m = kb \in \mathbb{Z}$.

Así, $ab$ es par.

\end{proof}

\newpage

\subsection{Demostración de que la suma de los primeros $n$ números impares es $n^2$}

\begin{proposition}
Para todo número natural $n$, se cumple que:
\[
1 + 3 + 5 + \cdots + (2n-1) = n^2
\]
\end{proposition}

\begin{proof}

Procedemos por inducción matemática.

\textbf{Caso base:} Para $n=1$:

\[
1 = 1^2
\]

Por lo tanto, la proposición se cumple para $n=1$.

\textbf{Paso inductivo:} Supongamos que para algún $k \in \mathbb{N}$ se cumple que:

\[
1 + 3 + 5 + \cdots + (2k-1) = k^2
\]

Queremos demostrar que también se cumple para $k+1$. El siguiente número impar después de $2k-1$ es:

\[
2(k+1)-1 = 2k+1
\]

Entonces:

\[
1 + 3 + 5 + \cdots + (2k-1) + (2k+1)
\]

Usando la hipótesis de inducción:

\[
= k^2 + 2k + 1
\]

Reconocemos el lado derecho como un cuadrado perfecto:

\[
k^2 + 2k + 1 = (k+1)^2
\]

Por lo tanto, la proposición se cumple para $k+1$.

Por el principio de inducción matemática, se concluye que:

\[
1 + 3 + 5 + \cdots + (2n-1) = n^2
\]

para todo $n \in \mathbb{N}$.

\end{proof}

\newpage

\subsection{Demostración de que la diferencia de dos números pares es par}

\begin{proposition}
La diferencia de dos números pares es par.
\end{proposition}

\begin{proof}

Sean $a$ y $b$ dos números pares. Por definición, existen enteros $k$ y $m$ tales que:

\[
a = 2k, \qquad b = 2m
\]

Entonces:

\[
a-b = 2k - 2m
\]

Factorizamos el $2$:

\[
a-b = 2(k-m)
\]

Como $k$ y $m$ son enteros, la diferencia $k-m$ también es un entero. Por lo tanto, $a-b$ tiene la forma $2q$, con $q \in \mathbb{Z}$.

Así, $a-b$ es par.

\end{proof}

\newpage

\subsection{Demostración de que el cuadrado de un número par es par}

\begin{proposition}
Si $n$ es un número par, entonces $n^2$ también es par.
\end{proposition}

\begin{proof}

Como $n$ es par, existe un entero $k$ tal que:

\[
n = 2k
\]

Elevamos ambos lados al cuadrado:

\[
n^2 = (2k)^2
\]

Desarrollamos:

\[
n^2 = 4k^2
\]

Reescribimos el resultado como múltiplo de $2$:

\[
n^2 = 2(2k^2)
\]

Como $2k^2$ es un entero, $n^2$ tiene la forma $2q$, con $q \in \mathbb{Z}$.

Por lo tanto, $n^2$ es par.

\end{proof}

\newpage

\subsection{Demostración de que si $n^2$ es par, entonces $n$ es par}

\begin{proposition}
Si $n$ es un número entero y $n^2$ es par, entonces $n$ es par.
\end{proposition}

\begin{proof}

Procedemos por contraposición.

La proposición contrapositiva dice: si $n$ no es par, entonces $n^2$ no es par. Es decir, si $n$ es impar, entonces $n^2$ es impar.

Supongamos que $n$ es impar. Entonces existe un entero $k$ tal que:

\[
n = 2k+1
\]

Elevamos al cuadrado:

\[
n^2 = (2k+1)^2
\]

Desarrollamos:

\[
n^2 = 4k^2 + 4k + 1
\]

Factorizamos los primeros dos términos:

\[
n^2 = 2(2k^2+2k)+1
\]

Como $2k^2+2k$ es un entero, $n^2$ tiene la forma $2m+1$. Por lo tanto, $n^2$ es impar.

Así, hemos demostrado que si $n$ es impar, entonces $n^2$ es impar. Por contraposición, si $n^2$ es par, entonces $n$ es par.

\end{proof}

\newpage

\subsection{Demostración de que la suma de un número par y un número impar es impar}

\begin{proposition}
La suma de un número par y un número impar es impar.
\end{proposition}

\begin{proof}

Sea $a$ un número par y sea $b$ un número impar.

Por definición, existen enteros $k$ y $m$ tales que:

\[
a = 2k, \qquad b = 2m+1
\]

Sumamos ambos números:

\[
a+b = 2k + (2m+1)
\]

Agrupamos términos semejantes:

\[
a+b = 2k + 2m + 1
\]

Factorizamos el $2$ en los dos primeros términos:

\[
a+b = 2(k+m)+1
\]

Como $k+m$ es un entero, $a+b$ tiene la forma $2q+1$, con $q \in \mathbb{Z}$.

Por lo tanto, $a+b$ es impar.

\end{proof}

\newpage

\chapter{Demostraciones algebraicas}

\section{Propiedades algebraicas}

\subsection{Demostración de la identidad $(a+b)^2 = a^2 + 2ab + b^2$}

\begin{proposition}
Para cualesquiera números reales $a$ y $b$, se cumple que:
\[
(a+b)^2 = a^2 + 2ab + b^2
\]
\end{proposition}

\begin{proof}

Partimos del lado izquierdo de la igualdad:

\[
(a+b)^2
\]

Por definición de potencia, elevar al cuadrado significa multiplicar la expresión por sí misma:

\[
(a+b)^2 = (a+b)(a+b)
\]

Ahora aplicamos la propiedad distributiva de la multiplicación respecto de la suma:

\[
(a+b)(a+b) = a(a+b) + b(a+b)
\]

Volvemos a aplicar la propiedad distributiva en cada término:

\[
a(a+b) = a^2 + ab
\]

\[
b(a+b) = ab + b^2
\]

Sumando ambos resultados:

\[
(a^2 + ab) + (ab + b^2)
\]

Agrupamos términos semejantes:

\[
a^2 + ab + ab + b^2
\]

\[
a^2 + 2ab + b^2
\]

Por lo tanto:

\[
(a+b)^2 = a^2 + 2ab + b^2
\]

\end{proof}

\newpage

\subsection{Demostración de que $a^2-b^2=(a-b)(a+b)$}

\begin{proposition}
Para cualesquiera números reales $a$ y $b$, se cumple que:
\[
a^2 - b^2 = (a-b)(a+b)
\]
\end{proposition}

\begin{proof}

Partimos del lado derecho de la igualdad:

\[
(a-b)(a+b)
\]

Aplicamos la propiedad distributiva, es decir, multiplicamos cada término del primer paréntesis por cada término del segundo:

\[
(a-b)(a+b) = a(a+b) - b(a+b)
\]

Desarrollamos cada uno de los productos:

\[
a(a+b) = a^2 + ab
\]

\[
b(a+b) = ab + b^2
\]

Sustituimos en la expresión original:

\[
(a-b)(a+b) = (a^2 + ab) - (ab + b^2)
\]

Distribuimos el signo negativo:

\[
= a^2 + ab - ab - b^2
\]

Ahora simplificamos términos semejantes:

\[
ab - ab = 0
\]

Por lo tanto:

\[
(a-b)(a+b) = a^2 - b^2
\]

Hemos demostrado que:

\[
a^2 - b^2 = (a-b)(a+b)
\]

\end{proof}

\newpage

\subsection{Demostración de la fórmula cuadrática}

\begin{proposition}
Las soluciones de la ecuación cuadrática
\[
ax^2 + bx + c = 0, \quad a \neq 0
\]
están dadas por:
\[
x = \frac{-b \pm \sqrt{b^2 - 4ac}}{2a}
\]
\end{proposition}


\begin{proof}

Partimos de la ecuación general:

\[
ax^2 + bx + c = 0
\]

Dividimos entre $a$:

\[
x^2 + \frac{b}{a}x + \frac{c}{a} = 0
\]

Pasamos el término constante al otro lado:

\[
x^2 + \frac{b}{a}x = -\frac{c}{a}
\]

Completamos el cuadrado. Tomamos la mitad del coeficiente de $x$ y la elevamos al cuadrado:

\[
\left(\frac{b}{2a}\right)^2 = \frac{b^2}{4a^2}
\]

Sumamos este término a ambos lados:

\[
x^2 + \frac{b}{a}x + \frac{b^2}{4a^2} = -\frac{c}{a} + \frac{b^2}{4a^2}
\]

El lado izquierdo es un cuadrado perfecto:

\[
\left(x + \frac{b}{2a}\right)^2 = \frac{b^2}{4a^2} - \frac{c}{a}
\]

Ponemos el lado derecho en común denominador:

\[
\left(x + \frac{b}{2a}\right)^2 = \frac{b^2 - 4ac}{4a^2}
\]

Sacamos raíz cuadrada:

\[
x + \frac{b}{2a} = \pm \frac{\sqrt{b^2 - 4ac}}{2a}
\]

Restamos $\frac{b}{2a}$:

\[
x = \frac{-b \pm \sqrt{b^2 - 4ac}}{2a}
\]

\end{proof}
\newpage

\subsection{Demostración de $a^3 + b^3 = (a+b)(a^2 - ab + b^2)$}

\begin{proposition}
Para cualesquiera números reales $a$ y $b$, se cumple que:
\[
a^3 + b^3 = (a+b)(a^2 - ab + b^2)
\]
\end{proposition}

\begin{proof}

Partimos del lado derecho de la igualdad:

\[
(a+b)(a^2 - ab + b^2)
\]

Aplicamos la propiedad distributiva:

\[
= a(a^2 - ab + b^2) + b(a^2 - ab + b^2)
\]

Distribuimos en cada término:

\[
= (a^3 - a^2b + ab^2) + (a^2b - ab^2 + b^3)
\]

Agrupamos términos semejantes:

\[
= a^3 - a^2b + ab^2 + a^2b - ab^2 + b^3
\]

Simplificamos:

\[
-a^2b + a^2b = 0, \qquad ab^2 - ab^2 = 0
\]

Por lo tanto:

\[
(a+b)(a^2 - ab + b^2) = a^3 + b^3
\]

\end{proof}

\newpage


\subsection{Demostración de $a^2+2ab+b^2=(a+b)^2$}

\begin{proposition}
Para cualesquiera números reales $a$ y $b$, se cumple que:
\[
a^2+2ab+b^2=(a+b)^2
\]
\end{proposition}

\begin{proof}

Partimos del lado izquierdo:

\[
a^2+2ab+b^2
\]

Escribimos $2ab$ como $ab+ab$:

\[
a^2+2ab+b^2 = a^2+ab+ab+b^2
\]

Agrupamos los primeros dos términos y los últimos dos términos:

\[
a^2+ab+ab+b^2 = (a^2+ab)+(ab+b^2)
\]

Factorizamos $a$ en el primer paréntesis y $b$ en el segundo:

\[
(a^2+ab)+(ab+b^2)=a(a+b)+b(a+b)
\]

Ahora factorizamos el factor común $(a+b)$:

\[
a(a+b)+b(a+b)=(a+b)(a+b)
\]

Por definición de potencia:

\[
(a+b)(a+b)=(a+b)^2
\]

Por lo tanto:

\[
a^2+2ab+b^2=(a+b)^2
\]

\end{proof}

\newpage

\subsection{Demostración de que $(a-b)^2=a^2-2ab+b^2$}

\begin{proposition}
Para cualesquiera números reales $a$ y $b$, se cumple que:
\[
(a-b)^2=a^2-2ab+b^2
\]
\end{proposition}

\begin{proof}

Por definición de potencia:

\[
(a-b)^2=(a-b)(a-b)
\]

Aplicamos la propiedad distributiva:

\[
(a-b)(a-b)=a(a-b)-b(a-b)
\]

Desarrollamos cada producto:

\[
a(a-b)=a^2-ab
\]

\[
b(a-b)=ab-b^2
\]

Sustituimos:

\[
(a-b)(a-b)=a^2-ab-(ab-b^2)
\]

Distribuimos el signo negativo:

\[
a^2-ab-ab+b^2
\]

Agrupamos términos semejantes:

\[
a^2-2ab+b^2
\]

Por lo tanto:

\[
(a-b)^2=a^2-2ab+b^2
\]

\end{proof}

\newpage

\chapter{Demostraciones Geométricas}
\section{Propiedades Geométricas}
\subsection{Teorema de Pitágoras}

\begin{proposition}
En un triángulo rectángulo, el cuadrado de la hipotenusa es igual a la suma de los cuadrados de los catetos.
\end{proposition}

\begin{proof}

Consideremos un cuadrado de lado $a+b$.

El área total del cuadrado es:
\[
(a+b)^2
\]

Ahora, dentro del cuadrado colocamos cuatro triángulos rectángulos congruentes, cada uno con catetos $a$ y $b$, e hipotenusa $c$.

El área de cada triángulo es:
\[
\frac{ab}{2}
\]

Entonces, el área de los cuatro triángulos es:
\[
4 \cdot \frac{ab}{2} = 2ab
\]

En el centro queda un cuadrado de lado $c$, cuya área es:
\[
c^2
\]

Por lo tanto, el área total también se puede expresar como:
\[
2ab + c^2
\]

Igualamos ambas expresiones del área:
\[
(a+b)^2 = 2ab + c^2
\]

Desarrollamos el lado izquierdo:
\[
a^2 + 2ab + b^2 = 2ab + c^2
\]

Restamos $2ab$ en ambos lados:
\[
a^2 + b^2 = c^2
\]

\end{proof}
\begin{figure}[h]
\centering
\includegraphics[width=0.6\textwidth]{pitagoras.png}
\caption{Representación geométrica del teorema de Pitágoras}
\end{figure}


\newpage

\subsection{Demostración del área de un triángulo}

\begin{proposition}
El área de un triángulo de base $b$ y altura $h$ está dada por:
\[
A = \frac{bh}{2}
\]
\end{proposition}

\begin{proof}

Consideremos un triángulo con base $b$ y altura $h$.

Si tomamos dos triángulos congruentes y los acomodamos adecuadamente, podemos formar un paralelogramo cuya base es $b$ y cuya altura es $h$.

El área de ese paralelogramo es:
\[
bh
\]

Como el paralelogramo está formado por dos triángulos congruentes, cada triángulo ocupa exactamente la mitad del área total.

Por lo tanto, el área de un triángulo es:
\[
A = \frac{bh}{2}
\]

\end{proof}

\newpage

\subsection{Demostración de la suma de los ángulos interiores de un triángulo}

\begin{proposition}
La suma de los ángulos interiores de todo triángulo es $180^\circ$.
\end{proposition}

\begin{proof}

Sea $ABC$ un triángulo cualquiera.

Trazamos por el vértice $A$ una recta paralela al lado $BC$.

Como esta recta es paralela a $BC$, los ángulos que se forman con los lados $AB$ y $AC$ son iguales a los ángulos interiores ubicados en $B$ y $C$, respectivamente, por ser ángulos alternos internos.

Sobre la recta que pasa por $A$, los tres ángulos consecutivos forman un ángulo llano. Por lo tanto, su suma es:
\[
180^\circ
\]

Pero esos tres ángulos corresponden exactamente a los tres ángulos interiores del triángulo.

Por lo tanto:
\[
\angle A + \angle B + \angle C = 180^\circ
\]

\end{proof}


\newpage

\chapter{Introducción a la física}

\section{Conceptos básicos de física}

\subsection{Relación entre distancia, velocidad y tiempo}

\begin{proposition}
Si un objeto se mueve con velocidad constante $v$ durante un tiempo $t$, entonces la distancia recorrida $d$ está dada por:
\[
d = vt
\]
\end{proposition}

\begin{proof}

La velocidad constante se define como la razón entre la distancia recorrida y el tiempo empleado:
\[
v = \frac{d}{t}
\]

Multiplicamos ambos lados por $t$:
\[
vt = d
\]

Por conmutatividad de la igualdad, obtenemos:
\[
d = vt
\]

Esta fórmula expresa que, si la velocidad no cambia, la distancia recorrida es proporcional al tiempo transcurrido.

\end{proof}

\newpage

\subsection{Segunda ley de Newton}

\begin{proposition}
La fuerza neta $F$ que actúa sobre un cuerpo de masa $m$ y aceleración $a$ está dada por:
\[
F = ma
\]
\end{proposition}

\begin{proof}

La segunda ley de Newton establece que la fuerza neta aplicada a un cuerpo es directamente proporcional a la aceleración que produce.

Si la masa del cuerpo es $m$, entonces una mayor masa requiere una mayor fuerza para producir la misma aceleración.

Por ello, la relación entre fuerza, masa y aceleración se expresa mediante:
\[
F = ma
\]

Esta ecuación indica que, para una masa fija, al aumentar la fuerza aumenta la aceleración; y para una aceleración fija, al aumentar la masa se necesita una fuerza mayor.

\end{proof}

\newpage

\subsection{Demostración de la fórmula del peso}

\begin{proposition}
El peso $P$ de un cuerpo de masa $m$ en un campo gravitatorio de aceleración $g$ está dado por:
\[
P = mg
\]
\end{proposition}

\begin{proof}

La segunda ley de Newton establece que la fuerza neta sobre un cuerpo se calcula mediante:
\[
F = ma
\]

El peso es la fuerza gravitatoria con la que un planeta atrae a un cuerpo.

En la superficie terrestre, la aceleración producida por la gravedad se representa por $g$.

Por lo tanto, en la fórmula $F=ma$ sustituimos la aceleración $a$ por la aceleración gravitatoria $g$:
\[
F = mg
\]

Como esta fuerza gravitatoria es precisamente el peso, escribimos:
\[
P = mg
\]

Así, el peso de un cuerpo depende de su masa y de la aceleración gravitatoria del lugar donde se encuentra.

\end{proof}




\newpage
\chapter*{Glosario}
\addcontentsline{toc}{chapter}{Glosario}

\section*{Conceptos de aritmética y teoría de números}

\textbf{Aritmética:} Rama de las matemáticas que estudia los números y las operaciones básicas entre ellos, como la suma, la resta, la multiplicación y la división. También analiza propiedades elementales de los números naturales, enteros, racionales y sus relaciones.

\vspace{0.5em}

\textbf{Álgebra:} Rama de las matemáticas que utiliza símbolos, letras y expresiones para representar números, cantidades desconocidas y relaciones entre ellas. Permite generalizar patrones, formular ecuaciones y demostrar identidades de manera ordenada.

\vspace{0.5em}

\textbf{Número par:} Número entero divisible entre $2$, es decir, aquel que puede escribirse en la forma $2k$ para algún entero $k$. Ejemplos de números pares son $0$, $2$, $-4$ y $10$.

\vspace{0.5em}

\textbf{Número impar:} Número entero que no es divisible entre $2$, o equivalentemente, número que puede escribirse en la forma $2k+1$ para algún entero $k$. Ejemplos de números impares son $1$, $3$, $-5$ y $11$.

\vspace{0.5em}

\textbf{Número racional:} Número que puede expresarse como una fracción $\frac{a}{b}$ con $a, b \in \mathbb{Z}$ y $b \neq 0$. Su desarrollo decimal es finito o periódico. Por ejemplo, $\frac{1}{2}$, $-3$ y $0.333\ldots$ son números racionales.

\vspace{0.5em}

\textbf{Número irracional:} Número que no puede expresarse como cociente de dos enteros. Su desarrollo decimal es infinito y no periódico. Ejemplos clásicos son $\sqrt{2}$, $\sqrt{3}$ y $\pi$.

\vspace{0.5em}

\textbf{Demostración:} Argumento lógico y ordenado que prueba la veracidad de una proposición matemática a partir de definiciones, axiomas, propiedades conocidas y resultados previamente demostrados. Su objetivo es justificar por qué una afirmación es verdadera.

\vspace{0.5em}

\textbf{Proposición:} Enunciado matemático con significado preciso que puede clasificarse como verdadero o falso. Muchas demostraciones comienzan estableciendo una proposición que luego será justificada paso a paso.

\vspace{0.5em}

\textbf{Teorema:} Proposición matemática importante que ha sido demostrada correctamente mediante un razonamiento lógico. Los teoremas suelen expresar resultados generales y sirven como base para obtener nuevas conclusiones.

\vspace{0.5em}

\textbf{Lema:} Resultado auxiliar que, aunque puede ser interesante por sí mismo, se utiliza principalmente para facilitar la demostración de un teorema u otra proposición más amplia.

\vspace{0.5em}

\textbf{Contradicción:} Método de demostración en el que se supone falsa la afirmación que se desea probar y, a partir de esa suposición, se llega a una imposibilidad lógica o a una afirmación falsa. Por ello, la proposición original debe ser verdadera.

\vspace{0.5em}

\textbf{Inducción matemática:} Método de demostración usado principalmente en los números naturales. Consiste en verificar primero un caso base y después demostrar que, si la afirmación se cumple para un número natural $n$, entonces también se cumple para $n+1$.

\vspace{0.5em}

\textbf{Propiedad distributiva:} Propiedad algebraica que establece que la multiplicación se distribuye sobre la suma, es decir, $a(b+c) = ab + ac$. También puede escribirse en la forma $(a+b)c = ac + bc$.

\vspace{0.5em}

\textbf{Factorización:} Proceso de reescribir una expresión algebraica como producto de factores más simples. Es una herramienta fundamental para simplificar expresiones, resolver ecuaciones y reconocer estructuras algebraicas importantes.

\vspace{0.5em}

\textbf{Cuadrado perfecto:} Número o expresión algebraica que puede escribirse como el cuadrado de otra cantidad. Por ejemplo, $9 = 3^2$ y $x^2 + 2xy + y^2 = (x+y)^2$.

\vspace{0.5em}

\textbf{Triángulo rectángulo:} Triángulo que tiene un ángulo de $90^\circ$, llamado ángulo recto. En este tipo de triángulo se aplica el teorema de Pitágoras para relacionar las longitudes de sus lados.

\vspace{0.5em}

\textbf{Hipotenusa:} Lado opuesto al ángulo recto en un triángulo rectángulo. Es el lado de mayor longitud y su cuadrado es igual a la suma de los cuadrados de los catetos.

\vspace{0.5em}

\textbf{Catetos:} Los dos lados de un triángulo rectángulo que forman el ángulo de $90^\circ$. En el teorema de Pitágoras, los cuadrados de sus longitudes suman el cuadrado de la hipotenusa.

\vspace{0.5em}

\textbf{Congruente:} Término que se usa para describir figuras geométricas que tienen exactamente la misma forma y el mismo tamaño, aunque puedan estar ubicadas en posiciones distintas o con diferente orientación.

\vspace{0.5em}

\textbf{Área:} Medida de la superficie ocupada por una figura plana. Dependiendo de la figura, puede calcularse mediante distintas fórmulas, como base por altura, productos notables o descomposición en partes más simples.

\vspace{0.5em}

\textbf{$\mathbb{Z}$:} Conjunto de los números enteros, es decir, los números negativos, el cero y los números positivos sin parte decimal. Se representa como $\mathbb{Z} = \{\ldots, -2, -1, 0, 1, 2, \ldots\}$.

\vspace{0.5em}

\textbf{$\mathbb{N}$:} Conjunto de los números naturales, usados principalmente para contar. Según la convención, puede incluir o no al $0$; en muchos contextos escolares se toma como $\mathbb{N} = \{1, 2, 3, \ldots\}$.

\vspace{0.5em}

\textbf{Número primo:} Número natural mayor que $1$ que tiene exactamente dos divisores positivos: el $1$ y él mismo. Ejemplos de números primos son $2$, $3$, $5$ y $7$.

\vspace{0.5em}

\textbf{Divisor:} Se dice que un número entero $d$ es divisor de otro entero $n$ si existe un entero $k$ tal que $n = dk$. En ese caso, también se dice que $d$ divide a $n$.

\vspace{0.5em}

\textbf{Múltiplo:} Un número es múltiplo de otro si puede escribirse como el producto de ese número por un entero. Por ejemplo, $12$ es múltiplo de $3$ porque $12 = 3 \cdot 4$.

\vspace{0.5em}

\textbf{Identidad algebraica:} Igualdad que es verdadera para todos los valores permitidos de sus variables. Por ejemplo, $(a+b)^2 = a^2 + 2ab + b^2$ es una identidad algebraica.

\vspace{0.5em}

\textbf{Ecuación:} Igualdad que contiene una o más incógnitas y que solo se cumple para ciertos valores de ellas. Resolver una ecuación consiste en encontrar todos los valores que hacen verdadera esa igualdad.

\vspace{0.5em}

\textbf{Conmutatividad:} Propiedad que indica que el resultado no cambia al alterar el orden de ciertos elementos. Por ejemplo, en los números reales, $a+b=b+a$ y $ab=ba$.

\vspace{0.5em}

\textbf{Asociatividad:} Propiedad que permite reagrupar sumas o productos sin cambiar el resultado. Por ejemplo, $(a+b)+c = a+(b+c)$ y $(ab)c = a(bc)$.
\vspace{0.5em}

\textbf{Variable:} Símbolo, generalmente una letra, que representa un número o una cantidad que puede cambiar. Por ejemplo, en la expresión $2x+1$, la letra $x$ es una variable.

\vspace{0.5em}

\textbf{Constante:} Valor que permanece fijo dentro de una expresión, ecuación o problema. Por ejemplo, en $3x+5$, el número $5$ es una constante.

\vspace{0.5em}

\textbf{Coeficiente:} Número que multiplica a una variable o a una expresión algebraica. En $7x$, el coeficiente de $x$ es $7$.

\vspace{0.5em}

\textbf{Términos semejantes:} Términos algebraicos que tienen la misma parte literal, es decir, las mismas variables elevadas a los mismos exponentes. Por ejemplo, $3x^2$ y $-5x^2$ son términos semejantes.

\vspace{0.5em}

\textbf{Potencia:} Expresión que representa una multiplicación repetida. Por ejemplo, $a^3$ significa $a \cdot a \cdot a$.

\vspace{0.5em}

\textbf{Exponente:} Número que indica cuántas veces se multiplica una base por sí misma. En $2^5$, el exponente es $5$.

\vspace{0.5em}

\textbf{Base:} Número o expresión que se repite en una potencia. En $x^4$, la base es $x$.

\vspace{0.5em}

\textbf{Raíz cuadrada:} Número no negativo que, al multiplicarse por sí mismo, produce el número dado. Por ejemplo, $\sqrt{9}=3$ porque $3^2=9$.

\vspace{0.5em}

\textbf{Desigualdad:} Relación matemática que compara dos cantidades mediante símbolos como $<$, $>$, $\leq$ o $\geq$.

\vspace{0.5em}

\textbf{Hipótesis:} Suposición o condición inicial que se acepta para desarrollar una demostración. En una proposición de la forma ``si $P$, entonces $Q$'', la hipótesis es $P$.

\vspace{0.5em}

\textbf{Conclusión:} Afirmación que se obtiene al final de un razonamiento o demostración. En una proposición de la forma ``si $P$, entonces $Q$'', la conclusión es $Q$.

\vspace{0.5em}

\textbf{Caso base:} Primer paso de una demostración por inducción matemática. Consiste en verificar que la proposición se cumple para el primer valor considerado.

\vspace{0.5em}

\textbf{Paso inductivo:} Parte de una demostración por inducción donde se prueba que, si la proposición se cumple para un valor $k$, entonces también se cumple para $k+1$.

\vspace{0.5em}

\textbf{Hipótesis de inducción:} Suposición temporal de que una proposición es verdadera para un valor $k$, usada para demostrar que también es verdadera para $k+1$.

\vspace{0.5em}

\textbf{Número compuesto:} Número natural mayor que $1$ que tiene más de dos divisores positivos. Por ejemplo, $4$, $6$, $8$ y $9$ son números compuestos.

\vspace{0.5em}

\textbf{Residuo:} Cantidad que sobra después de realizar una división entera. Por ejemplo, al dividir $17$ entre $5$, el residuo es $2$.

\vspace{0.5em}

\textbf{Módulo:} Forma de estudiar los residuos de una división. La expresión $a \equiv b \pmod{n}$ significa que $a$ y $b$ dejan el mismo residuo al dividirse entre $n$.

\vspace{0.5em}

\textbf{Fracción irreducible:} Fracción cuyo numerador y denominador no tienen divisores comunes distintos de $1$. Por ejemplo, $\frac{2}{3}$ es irreducible.

\vspace{0.5em}

\textbf{Máximo común divisor:} Mayor número entero positivo que divide exactamente a dos o más números enteros. Se suele abreviar como MCD.

\vspace{0.5em}

\textbf{Mínimo común múltiplo:} Menor número entero positivo que es múltiplo común de dos o más números enteros. Se suele abreviar como MCM.

\vspace{0.5em}

\textbf{Coprimos:} Dos números enteros son coprimos si su máximo común divisor es $1$.

\vspace{0.5em}

\textbf{Paridad:} Propiedad de un número entero de ser par o impar.

\vspace{0.5em}

\textbf{Contraposición:} Método lógico que demuestra una proposición ``si $P$, entonces $Q$'' probando la proposición equivalente ``si no $Q$, entonces no $P$''.

\section*{Conceptos de álgebra}

\vspace{0.5em}

\textbf{Valor absoluto:} Distancia de un número al cero en la recta numérica. Se denota por $|a|$ y siempre es mayor o igual que cero.

\vspace{0.5em}

\textbf{Recta numérica:} Representación geométrica de los números sobre una línea recta, donde cada punto corresponde a un número.

\vspace{0.5em}

\textbf{Expresión algebraica:} Combinación de números, variables y operaciones matemáticas. Por ejemplo, $3x^2-5x+1$ es una expresión algebraica.

\vspace{0.5em}

\textbf{Polinomio:} Expresión algebraica formada por la suma o resta de términos con variables elevadas a exponentes enteros no negativos. Por ejemplo, $x^2+3x+2$ es un polinomio.

\vspace{0.5em}

\textbf{Monomio:} Expresión algebraica que consta de un solo término. Por ejemplo, $5x^3$ es un monomio.

\vspace{0.5em}

\textbf{Binomio:} Expresión algebraica formada por dos términos. Por ejemplo, $a+b$ y $x^2-1$ son binomios.

\vspace{0.5em}

\textbf{Trinomio:} Expresión algebraica formada por tres términos. Por ejemplo, $x^2+2x+1$ es un trinomio.

\vspace{0.5em}

\textbf{Producto notable:} Identidad algebraica frecuente que permite desarrollar o factorizar expresiones de forma rápida, como $(a+b)^2=a^2+2ab+b^2$.

\vspace{0.5em}

\textbf{Término independiente:} Término de una expresión algebraica que no contiene variables. En $x^2+3x+4$, el término independiente es $4$.

\vspace{0.5em}

\textbf{Grado de un polinomio:} Mayor exponente de la variable en un polinomio de una variable. Por ejemplo, el grado de $2x^3+x-1$ es $3$.

\vspace{0.5em}

\textbf{Solución de una ecuación:} Valor que, al sustituirse en una ecuación, hace que la igualdad sea verdadera.

\vspace{0.5em}

\textbf{Sistema de ecuaciones:} Conjunto de dos o más ecuaciones que se estudian al mismo tiempo. Resolverlo consiste en encontrar valores que satisfagan todas las ecuaciones del sistema.

\vspace{0.5em}

\textbf{Ecuación lineal:} Ecuación en la que la incógnita aparece elevada únicamente a la primera potencia. Por ejemplo, $2x+3=7$.

\vspace{0.5em}

\textbf{Ecuación cuadrática:} Ecuación que puede escribirse en la forma $ax^2+bx+c=0$, con $a \neq 0$.

\vspace{0.5em}

\textbf{Discriminante:} En una ecuación cuadrática, es la expresión $b^2-4ac$. Permite conocer la naturaleza de sus soluciones.

\vspace{0.5em}

\textbf{Incógnita:} Cantidad desconocida que se busca determinar al resolver una ecuación.

\vspace{0.5em}

\textbf{Simplificación:} Proceso de transformar una expresión en otra equivalente pero más sencilla.

\vspace{0.5em}

\textbf{Equivalencia:} Relación entre dos expresiones o proposiciones que tienen el mismo valor o significado matemático.

\section*{Conceptos de funciones y geometría analítica}

\vspace{0.5em}

\textbf{Función:} Relación que asigna a cada elemento de un conjunto de entrada exactamente un elemento de un conjunto de salida.

\vspace{0.5em}

\textbf{Dominio:} Conjunto de valores para los cuales una función o expresión está definida.

\vspace{0.5em}

\textbf{Rango:} Conjunto de valores que puede tomar una función como resultado.

\vspace{0.5em}

\textbf{Plano cartesiano:} Sistema formado por dos rectas numéricas perpendiculares, llamadas eje $x$ y eje $y$, que permite representar puntos mediante coordenadas.

\vspace{0.5em}

\textbf{Coordenada:} Número que indica la posición de un punto respecto de un eje. En el plano cartesiano, un punto suele escribirse como $(x,y)$.

\vspace{0.5em}

\textbf{Pendiente:} Medida de la inclinación de una recta. Indica cuánto cambia la coordenada vertical cuando cambia la coordenada horizontal.

\vspace{0.5em}

\textbf{Intersección:} Punto en el que dos figuras, rectas o curvas se encuentran.

\vspace{0.5em}

\textbf{Eje $x$:} Recta horizontal del plano cartesiano.

\vspace{0.5em}

\textbf{Eje $y$:} Recta vertical del plano cartesiano.

\vspace{0.5em}

\textbf{Origen:} Punto de intersección entre el eje $x$ y el eje $y$ en el plano cartesiano. Se representa como $(0,0)$.

\vspace{0.5em}

\textbf{Ordenada al origen:} Punto donde una recta corta al eje $y$.

\section*{Conceptos de conjuntos y lógica}

\vspace{0.5em}

\textbf{Paralelismo:} Relación entre rectas que están en el mismo plano y nunca se cortan.

\vspace{0.5em}

\textbf{Perpendicularidad:} Relación entre dos rectas que se cortan formando un ángulo recto de $90^\circ$.

\vspace{0.5em}

\textbf{Conjunto:} Colección de objetos matemáticos bien definidos, llamados elementos. Por ejemplo, $\{1,2,3\}$ es un conjunto.

\vspace{0.5em}

\textbf{Elemento:} Objeto que pertenece a un conjunto. Si $3$ pertenece al conjunto $A$, se escribe $3 \in A$.

\vspace{0.5em}

\textbf{Subconjunto:} Conjunto cuyos elementos pertenecen todos a otro conjunto. Si todo elemento de $A$ pertenece a $B$, entonces $A$ es subconjunto de $B$.

\vspace{0.5em}

\textbf{Unión:} Operación entre conjuntos que reúne los elementos que pertenecen a uno, al otro o a ambos conjuntos.

\vspace{0.5em}

\textbf{Intersección de conjuntos:} Operación que reúne únicamente los elementos comunes a dos conjuntos.

\vspace{0.5em}

\textbf{Axioma:} Afirmación básica que se acepta como verdadera dentro de un sistema matemático y sirve como punto de partida para demostrar otros resultados.

\vspace{0.5em}

\textbf{Corolario:} Resultado que se obtiene fácilmente a partir de un teorema o proposición ya demostrada.

\vspace{0.5em}

\textbf{Recíproca:} Proposición que se obtiene al intercambiar la hipótesis y la conclusión de una implicación. La recíproca de ``si $P$, entonces $Q$'' es ``si $Q$, entonces $P$''.

\vspace{0.5em}

\textbf{Contraejemplo:} Ejemplo específico que demuestra que una afirmación general es falsa.

\vspace{0.5em}

\textbf{Implicación:} Proposición lógica de la forma ``si $P$, entonces $Q$''. Indica que la verdad de $P$ conduce a la verdad de $Q$.

\vspace{0.5em}

\textbf{Bicondicional:} Proposición de la forma ``$P$ si y solo si $Q$''. Significa que $P$ implica $Q$ y $Q$ implica $P$.

\vspace{0.5em}

\textbf{Necesario:} Una condición necesaria es aquella que debe cumplirse para que una afirmación sea verdadera.

\vspace{0.5em}

\textbf{Suficiente:} Una condición suficiente es aquella que garantiza que una afirmación sea verdadera.

\section*{Conceptos de geometría}

\vspace{0.5em}

\textbf{Base:} En geometría, lado de una figura que se toma como referencia para calcular áreas o alturas. En un triángulo, cualquier lado puede elegirse como base.

\vspace{0.5em}

\textbf{Altura:} Segmento perpendicular que va desde un vértice hasta la recta que contiene la base opuesta. Es fundamental para calcular áreas de triángulos y paralelogramos.

\vspace{0.5em}

\textbf{Paralelogramo:} Cuadrilátero cuyos lados opuestos son paralelos. Su área se calcula multiplicando la base por la altura.

\vspace{0.5em}

\textbf{Ángulo:} Abertura formada por dos semirrectas que comparten un mismo origen, llamado vértice.

\vspace{0.5em}

\textbf{Ángulo llano:} Ángulo que mide $180^\circ$ y forma una línea recta.

\vspace{0.5em}

\textbf{Ángulos alternos internos:} Pares de ángulos formados cuando una recta corta a dos rectas paralelas. Estos ángulos tienen la misma medida.

\vspace{0.5em}

\textbf{Vértice:} Punto donde se encuentran dos lados de una figura geométrica o dos semirrectas que forman un ángulo.

\vspace{0.5em}

\textbf{Lado:} Segmento que forma parte del contorno de una figura geométrica.

\vspace{0.5em}

\textbf{Cuadrilátero:} Polígono de cuatro lados. Ejemplos de cuadriláteros son el cuadrado, el rectángulo, el rombo y el paralelogramo.

\vspace{0.5em}

\textbf{Polígono:} Figura plana cerrada formada por segmentos de recta.

\vspace{0.5em}

\textbf{Circunferencia:} Conjunto de puntos del plano que están a la misma distancia de un punto fijo llamado centro.

\vspace{0.5em}

\textbf{Círculo:} Región del plano limitada por una circunferencia.

\vspace{0.5em}

\textbf{Radio:} Segmento que une el centro de una circunferencia con cualquiera de sus puntos.

\vspace{0.5em}

\textbf{Diámetro:} Segmento que une dos puntos de una circunferencia pasando por su centro. Su longitud es el doble del radio.

\vspace{0.5em}

\textbf{Perímetro:} Longitud total del contorno de una figura plana.

\section*{Conceptos de física}

\vspace{0.5em}

\textbf{Física:} Ciencia que estudia la materia, la energía, el movimiento, las fuerzas y las leyes que describen el comportamiento de la naturaleza.

\vspace{0.5em}

\textbf{Magnitud física:} Propiedad que puede medirse y expresarse mediante un número y una unidad, como la masa, el tiempo, la distancia o la velocidad.

\vspace{0.5em}

\textbf{Unidad de medida:} Cantidad estándar usada para medir una magnitud física. Por ejemplo, el metro mide longitud y el segundo mide tiempo.

\vspace{0.5em}

\textbf{Distancia:} Longitud total recorrida por un objeto durante su movimiento. Generalmente se mide en metros.

\vspace{0.5em}

\textbf{Tiempo:} Magnitud física que permite ordenar sucesos y medir la duración de un fenómeno. Su unidad básica en el Sistema Internacional es el segundo.

\vspace{0.5em}

\textbf{Velocidad:} Magnitud que indica qué tan rápido cambia la posición de un objeto. En movimiento uniforme puede calcularse como $v=\frac{d}{t}$.

\vspace{0.5em}

\textbf{Rapidez:} Magnitud escalar que mide la distancia recorrida por unidad de tiempo, sin considerar la dirección del movimiento.

\vspace{0.5em}

\textbf{Aceleración:} Magnitud que mide el cambio de velocidad respecto al tiempo. Si la velocidad cambia rápidamente, la aceleración es mayor.

\vspace{0.5em}

\textbf{Masa:} Cantidad de materia de un cuerpo. En el Sistema Internacional se mide en kilogramos.

\vspace{0.5em}

\textbf{Fuerza:} Interacción capaz de modificar el estado de movimiento de un cuerpo o deformarlo. En el Sistema Internacional se mide en newtons.

\vspace{0.5em}

\textbf{Newton:} Unidad de fuerza en el Sistema Internacional. Un newton es la fuerza necesaria para producir una aceleración de $1\,\mathrm{m/s^2}$ sobre una masa de $1\,\mathrm{kg}$.

\vspace{0.5em}

\textbf{Movimiento uniforme:} Movimiento en el que un objeto recorre distancias iguales en tiempos iguales, manteniendo velocidad constante.

\vspace{0.5em}

\textbf{Movimiento rectilíneo:} Movimiento cuya trayectoria es una línea recta.

\vspace{0.5em}

\textbf{Trayectoria:} Camino seguido por un objeto durante su movimiento.

\vspace{0.5em}

\textbf{Posición:} Lugar que ocupa un objeto respecto de un sistema de referencia.

\vspace{0.5em}

\textbf{Sistema de referencia:} Conjunto de criterios o ejes desde los cuales se describe la posición y el movimiento de un objeto.

\vspace{0.5em}

\textbf{Energía:} Capacidad de un cuerpo o sistema para realizar trabajo o producir cambios.

\vspace{0.5em}

\textbf{Trabajo:} En física, cantidad de energía transferida cuando una fuerza desplaza un cuerpo. Se calcula como $W=Fd$ cuando la fuerza y el desplazamiento tienen la misma dirección.

\vspace{0.5em}

\textbf{Peso:} Fuerza con la que la gravedad atrae a un cuerpo. Depende de la masa y de la aceleración gravitatoria.

\vspace{0.5em}

\textbf{Gravedad:} Interacción por la cual los cuerpos con masa se atraen entre sí. En la Tierra produce la caída de los objetos hacia el suelo.

\vspace{0.5em}

\textbf{Sistema Internacional de Unidades:} Sistema de unidades usado internacionalmente para expresar mediciones físicas de forma estandarizada.

\vspace{0.5em}

\textbf{Desplazamiento:} Cambio de posición de un objeto, considerando dirección y sentido.

\vspace{0.5em}

\textbf{Magnitud escalar:} Magnitud física que se describe solo con un número y una unidad, como la masa o el tiempo.

\vspace{0.5em}

\textbf{Magnitud vectorial:} Magnitud física que requiere número, unidad, dirección y sentido, como la fuerza o la velocidad.

\vspace{0.5em}

\textbf{Campo gravitatorio:} Región del espacio donde un cuerpo con masa experimenta una fuerza gravitatoria.

\vspace{0.5em}

\textbf{Aceleración gravitatoria:} Aceleración producida por la gravedad. Cerca de la superficie terrestre suele aproximarse por $9.8\,\mathrm{m/s^2}$.

\vspace{0.5em}

\textbf{Cantidad de movimiento:} Magnitud física definida como el producto de la masa por la velocidad de un cuerpo.

\vspace{0.5em}

\textbf{Inercia:} Tendencia de un cuerpo a conservar su estado de reposo o movimiento rectilíneo uniforme.

\vspace{0.5em}

\textbf{Potencia:} En física, rapidez con la que se realiza trabajo o se transfiere energía.


% LICENCIA 

\chapter*{Licencia}
\addcontentsline{toc}{chapter}{Licencia}

© 2026 Álvaro Yamil Manzo Ayora

Este documento está bajo la licencia

\textit{Creative Commons Attribution 4.0 International (CC BY 4.0)}.

Usted es libre de compartir y adaptar este material siempre que se otorgue el crédito correspondiente al autor.

Para ver una copia de esta licencia visite:

\begin{center}
\url{https://creativecommons.org/licenses/by/4.0/}
\end{center}

\end{document}
